\subsection{Module D.3: nonsupercritical high-radial output}
\label{subsec:app-finite-output-functions}

\paragraph{Exact interface.}
Use the canonical admissible set $\mathcal A$ and actual reaches of
\zcref{def:finite-anchor-admissible-set}.  The output function below is the
attained maximum outgoing reach on its nonsupercritical fiber.

\paragraph{Variables and feasible set.}
Let $a,b,c\in[0,1]$ satisfy $(a,b,c)\in\mathcal A$ and $a+b\le1$, with
$c>1/2$ for the high-radial branches.  In the threshold form put $c=1-e$
with $0<e<1-\sqrt3/2$; closed equality is allowed in the capacity, whereas
the hypothesis on the actual incoming reach is strict.

For a nonsupercritical V triangle with incoming reach at least $a$ and own
radial reach at least $c$, define the attained outgoing capacity
\[
 \mathcal B_c(a):=
 \max\{b\in[0,1]:(a,b,c)\in\mathcal A\ \land\ a+b\le1\}.
\]
For $0<e<1-\sqrt3/2$, put
\[
 \epsilon(e):=\frac{1-e}{2}
 \left(1-\sqrt{4(1-e)^2-3}\right).
\]
Set
\[
 \mu:=\frac{2\sqrt3-3}{4}.
\]
The selected low root obeys
\[
 e<\epsilon(e),
\]
\[
 0<e\le\mu\Longrightarrow
 \epsilon(e)<\frac{2e}{1-2e},
 \qquad
 0<e\le\frac18\Longrightarrow\epsilon(e)\le2e+5e^2.
\]

For $1/2<c\le\sqrt3/2$, put
\[
 h(c):=\frac{2c}{1+\sqrt{16c^2-3}},
\]
and, for $\sqrt3/2\le c<1$, put
\[
 \ell(c):=\frac c2\left(1-\sqrt{4c^2-3}\right),
 \qquad r(c):=c-\ell(c).
\]
Finally set
\[
 Q_-(a,c):=\frac{\sqrt{a^2-ac+c^2}}c-a,
\]
\[
 Q_+(a,c):=
 \frac{2(1-a^2)c^2}
 {2ac^2+c+
 \sqrt{(2ac^2+c)^2-4(1-c^2)(1-a^2)c^2}}.
\]

\paragraph{Objective.}
Maximize the outgoing reach $b$ on each high-radial nonsupercritical fiber,
then determine the exact low-root threshold that forces the outgoing actual
reach below the same root.

\paragraph{Exact result.}

\begin{proposition}[Four direct high-radial outputs]
\label{prop:app-four-direct-outputs}
For $1/2<c\le\sqrt3/2$,
\[
 \mathcal B_c(a)=
 \begin{cases}
  1-a,&0\le a\le1-c,\\
  Q_+(a,c),&1-c<a<h(c),\\
  h(c),&a=h(c),\\
  Q_-(a,c),&h(c)<a<c,\\
  1-a,&c\le a\le1.
 \end{cases}
\]
For $\sqrt3/2<c<1$,
\[
 \mathcal B_c(a)=
 \begin{cases}
  1-a,&0\le a\le1-c,\\
  Q_+(a,c),&1-c<a\le\ell(c),\\
  \ell(c),&\ell(c)<a<r(c),\\
  Q_-(a,c),&r(c)\le a<c,\\
  1-a,&c\le a\le1.
 \end{cases}
\]
At $a=\ell(c)$ the selected value is $r(c)$; immediately to its right it is
$\ell(c)$.
\end{proposition}

\begin{lemma}[Direct high-radial threshold]
\label{lem:app-direct-threshold}
Let a nonsupercritical V triangle have incoming, outgoing, and own-radial
reaches $A,B,C$.  If
\[
 0<e<1-\frac{\sqrt3}{2},
 \qquad C\ge1-e,
 \qquad A>\epsilon(e),
\]
then
\[
 B\le\epsilon(e).
\]
\end{lemma}

\paragraph{Solution.}

\begin{proof}[Proof of \zcref{prop:app-four-direct-outputs}]
Because $c>1/2$, the radial midpoint is contained and hence $a+b\le1$.  On
the cap $b=1-a$, the $T$-cell inequality from
\zcref{lem:app-local-admissible-set} reduces, with
$M=\max\{a,b\}$, to $M(c-M)\le0$.  Thus the cap is feasible exactly for
$a\le1-c$ or $a\ge c$.

Off that cap, a maximal point lies on $F_L=0$ or $F_T=0$.  When
$1/2<c<\sqrt3/2$, the discriminant $c^2(4c^2-3)$ of the $L$ contact is
negative, so no $L$-cell root occurs; the two ordered halves of $F_T=0$ give
$Q_+$ and $Q_-$ and meet at $h(c)$.  When $c\ge\sqrt3/2$, the selected
$L$ cell has $m=\min\{a,b\}\le\ell(c)$, and its selected radial frontier is
only $m=\ell(c)$.  Indeed, the algebraic $L$ roots are
$\ell(c),r(c)$, but the formal alternative $m\ge r(c)$ is
incompatible with the selector $q\le0$, except at the coalescence point.
The two ordered $q=0$ contacts are
$(\ell(c),r(c))$ and $(r(c),\ell(c))$, producing the exceptional left
endpoint and the constant interval.

In the ordered half $a\le b$, the cap $b\le1-a$ lies strictly between the
two formal roots of the $F_T$ quadratic, so its lower root is the selected
$Q_+$ branch and the upper root is excluded.  The remaining geometric
selector $c\le2\max\{a,b\}=2b$ holds throughout that retained branch: in the
first regime $b\ge h(c)\ge c/2$, and in the second
$b\ge r(c)\ge c/2$, with the latter equality only at
$c=\sqrt3/2$.  Reflection gives the selected $Q_-$ branch, for which
$c\le2a$.  These selectors give exactly the displayed pieces, including the
stated downward jump.
\end{proof}

\begin{proof}[Proof of \zcref{lem:app-direct-threshold}]
The value $\epsilon(e)=\ell(1-e)$ is the smaller root of
\[
 p_e(z):=z^2-(1-e)z+(1-e)^2-(1-e)^4=0.
\]
Substitution at $e$ gives
\[
 p_e(e)=e(1-3e+4e^2-e^3)>0,
\]
and $e<(1-e)/2$, the midpoint of the two roots.  Hence
$e<\epsilon(e)$ throughout the stated range.  If $0<e\le\mu$, then
\[
 p_e\!\left(\frac{2e}{1-2e}\right)
 =-\frac{e^2}{(1-2e)^2}
 P(e),
 \qquad
 P(e):=4e^4-20e^3+37e^2-28e+3.
\]
On $[0,\mu]\subset[0,1/8]$,
\[
 P'(e)=16e^3-60e^2+74e-28
 <\frac{16}{512}+\frac{74}{8}-28<0,
\]
whereas
\[
 P(\mu)=\frac{8217}{64}-74\sqrt3>0,
 \qquad
 \left(\frac{8217}{64}\right)^2-3\cdot74^2
 =\frac{230001}{4096}>0.
\]
Thus the preceding value of $p_e$ is strictly negative,
which proves the rational upper bound on precisely that range.  If
$e\le1/8$, then
\[
 p_e(2e+5e^2)=e^2(3e+4)(8e-1)\le0,
\]
proving the third.  Proposition~\ref{prop:app-four-direct-outputs}, with
$c=1-e$, shows that after $A>\epsilon(e)$ the selected constant and $Q_-$
pieces have outgoing coordinate at most $\epsilon(e)$.  The $Q_+$ piece has
already ended.  The upper linear branch $A\ge c$ remains possible, but there
$B\le1-A\le e<\epsilon(e)$.  Downward closure handles $C\ge1-e$.
\end{proof}

\paragraph{Body consequence.}
The evaluated output catalogue and threshold supply the high-radial steps in
\zcref{prop:readable-k410-ce2,prop:readable-k410-ce1}; later Appendix D
modules cite these two exact results rather than re-evaluating the fiber.

\subsection{Module D.4: strict-supercritical envelope}

\paragraph{Exact interface.}
The body defines $M_c^{\mathrm{sup}}$ as the supremum of outgoing reaches on
the strict-supercritical part of $\mathcal A$ in the supported-tail interface
preceding \zcref{prop:new-one-t3-terminal}.

\paragraph{Variables and feasible set.}
Take $0\le c<1/2$ and pairs $(a,b)\in[0,1]^2$ satisfying
$(a,b,c)\in\mathcal A$ and $a+b>1$.  For the relevant boundary reaches
$1/2\le b<1$, put
\[
 \kappa(b):=\frac{1-b^2}{2-b}.
\]
The support-slack feasibility condition is the strict inequality
$c<\kappa(b)$; it is not replaced by a blanket closure assertion for the
whole face $a+b=1$.  At $c=1/2$ use the body's declared continuous endpoint
value.

\paragraph{Objective.}
Determine the supremal outgoing reach $b$, decide attainment on the strict
fiber, and evaluate its one-sided endpoint limit.

\paragraph{Exact result.}

\begin{lemma}[Strict-supercritical envelope]
\label{lem:app-supercritical-envelope}
For $0\le c<1/2$ the variational quantity defined in the body is
\[
 M_c^{\mathrm{sup}}=\frac{c+\sqrt{c^2-8c+4}}2.
\]
It is not attained by a strict-supercritical pair.  Moreover
\[
 \lim_{c\uparrow1/2}M_c^{\mathrm{sup}}=\frac12
 =M_{1/2}^{\mathrm{sup}}.
\]
Thus a strict-supercritical V triangle with own-radial reach at least $c$
has outgoing boundary reach $B<M_c^{\mathrm{sup}}$.
\end{lemma}

\paragraph{Solution.}

\begin{proof}[Proof of \zcref{lem:app-supercritical-envelope}]
For $c=0$, downward closure and the positive-$c$ calculation below give
strict-supercritical outgoing reaches arbitrarily close to $1$, while always
$b\le1$.  This supremum is not attained: $b=1$ and
$a^2+ab+b^2\le1$ force $a=0$, contrary to $a+b>1$.
Fix $0<c<1/2$.  Only $b\ge1/2$ matters for the supremum.  Put
$a=1-b$, where $1/2\le b<1$, and use the support gauge of
\zcref{lem:app-shared-enclosure-gauge}.

Let $e_A,e_B$ be the two unit boundary directions from the corner, at angle
$2\pi/3$, so the boundary anchors are $ae_A$ and $be_B$.  We first prove
that $(a,b,c)$ has strict support slack, and therefore remains
feasible after increasing $a$ slightly into $a+b>1$, exactly when
$c<\kappa(b)$.  Reduce a support normal modulo $2\pi/3$ to the two sectors
$[-\pi/6,\pi/6]$ and $[\pi/6,\pi/2]$.  On the first sector the three-point
support sum for $0,ae_A,be_B$ is $\cos\phi\ge\sqrt3/2$.  On the second, put
\[
 p=\cos\phi,
 \qquad q=\cos(2\pi/3-\phi),
 \qquad t=\frac qp,
 \qquad p^2+pq+q^2=\frac34.
\]
After division by $p$, the four-point support sum is
\[
 N(t)=a+bt+\max\{at,b,c(1+t)\}.
\]
Set
\[
 t_*:=\frac{2b-1}{1-b^2}.
\]
When the maximum in $N$ is $b$,
\[
 (1+bt)^2-(1+t+t^2)
 =t\bigl((2b-1)-(1-b^2)t\bigr).
\]
Thus strict slack is possible for $t>t_*$.  The two inactive-support
conditions $at\le b$ and $c(1+t)\le b$ admit such a $t$ exactly when
\[
 t_*<\frac bc-1
 \quad\Longleftrightarrow\quad
 c<\kappa(b);
\]
the bound $t_*<b/a$ is automatic.  Continuity then permits a small strict
increase of $a$.

Conversely, suppose $c\ge\kappa(b)$.  For $t\le t_*$ the term $b$ gives
$N(t)\ge\sqrt{1+t+t^2}$.  For $t\ge t_*$, monotonicity reduces to
$c=\kappa(b)$.  With
\[
 A=a+\kappa(b),
 \qquad B=b+\kappa(b),
\]
the quadratic
\[
 (A+Bt)^2-(1+t+t^2)
\]
has negative constant term, positive leading coefficient, and positive root
$t=t_*$, because $\kappa(b)(1+t_*)=b$.  It is therefore nonnegative for
$t\ge t_*$.  Increasing $a$ makes the first-sector support sum strict and
strictly increases every second-sector expression (the endpoint $p=0$ is
already strict because $b+c>1$).  Hence no strict-supercritical triple with
this $b$ is feasible.

The decreasing function $\kappa$ maps $[1/2,1)$ onto $(0,1/2]$.
Consequently the supremum is the larger root of
\[
 b^2-cb+2c-1=0.
\]
Equality has no strict support slack, so this value is not attained.  The
limit at $c=1/2$ follows by direct substitution and equals the body's
declared continuous value.
\end{proof}

\paragraph{Body consequence.}
This formula and its strict nonattainment supply the value $M$ and the bound
$B_1<M$ in \zcref{prop:new-one-t3-terminal,prop:finite-vd1-supported-terminal}.

\subsection{Module D.5: selected-arc chord bounds}

\paragraph{Exact interface.}
Use the selected $Q_+$ component of the nonsupercritical output function from
Module D.3 and the center-free boundary handoffs in
\zcref{prop:readable-k410-ce1}.

\paragraph{Variables and feasible set.}
Let $0<e<1/2$, take $x\in[0,1]$, and set
$q=e+(1-e)x$.  The incoming coordinate $p$ is the smaller-root point on the
selected $Q_+$ component.  For $e<1-\sqrt3/2$ it is the high-radial arc; for
$1-\sqrt3/2\le e<1/4$ the CE1 use is the low-radial arc.  These ranges keep
every occurrence of $\epsilon(e)$ inside its declared domain.

\paragraph{Objective.}
Parametrize the selected output arc and obtain affine lower bounds for its
outgoing reach along the two intervals used by the CE1 reverse path.

\paragraph{Exact result.}

\begin{lemma}[Selected-arc chord bounds]
\label{lem:app-selected-chords}
Let $0<e<1/2$.  On the selected $Q_+$ component for radial demand $1-e$,
let $p$ be the incoming reach and let $q$ be the next incoming reach forced
across a center-free edge.  There is $x\in[0,1]$ with
\[
 q=e+(1-e)x,
 \qquad
 p=e+(1-e)x-1+\sqrt{1-x+x^2},
\]
and hence
\[
 q-p=1-\sqrt{1-x+x^2}.
\]
This selected arc is increasing and concave.  On each interval used by the
CE1 reverse path,
\[
 q\ge p+\lambda(p-e),
\]
with the following fully quantified choices.  On the high-radial arc
$0<e<1-\sqrt3/2$, take
\[
 \lambda_{\mathrm{hi}}(e):=\frac{e}{\epsilon(e)-e};
\]
then
\[
 0<e\le\mu\Longrightarrow\lambda_{\mathrm{hi}}(e)>1-4e,
 \qquad
 0<e<1-\frac{\sqrt3}{2}
 \Longrightarrow\lambda_{\mathrm{hi}}(e)>1-5e.
\]
On the low-radial arc $1-\sqrt3/2\le e<1/4$, put
\[
 t_e:=h(1-e),
 \qquad
 \lambda_{\mathrm{lo}}(e):=\frac{1-2t_e}{t_e-e};
\]
then
\[
 \lambda_{\mathrm{lo}}(e)>\frac13>1-5e.
\]
\end{lemma}

\paragraph{Solution.}

\begin{proof}[Proof of \zcref{lem:app-selected-chords}]
With $x=(q-e)/(1-e)$ the selected $T$-cell equation is
\[
 x(1-x)=(q-p)(2-q+p).
\]
The smaller root gives the displayed parametrization.  If
$\psi(x)=1-\sqrt{1-x+x^2}$, then
\[
 \psi''(x)=-\frac3{4(1-x+x^2)^{3/2}}<0.
\]
The map $x\mapsto p=e+(1-e)x-\psi(x)$ is strictly increasing and convex;
its inverse, and therefore $p\mapsto q$, is strictly increasing and concave.
The graph lies above each endpoint chord.

On the high-radial arc the endpoint pairs are
\[
 (p,q)=(e,e),
 \qquad (p,q)=(\epsilon(e),e+\epsilon(e)),
\]
which gives $\lambda_{\mathrm{hi}}(e)$.  For $e\le\mu$, the rational
low-root bound proved in Module D.3 gives
\[
 \lambda_{\mathrm{hi}}(e)
 >\frac{1-2e}{1+2e}>1-4e.
\]
If $e\le1/8$, the bound
$\epsilon(e)\le2e+5e^2$ gives
\[
 \lambda_{\mathrm{hi}}(e)\ge\frac1{1+5e}>1-5e.
\]
If $1/8<e<1-\sqrt3/2$, then
$\epsilon(e)<(1-e)/2$ directly from its defining radical, and hence
\[
 \lambda_{\mathrm{hi}}(e)>
 \frac{2e}{1-3e}>1-5e.
\]

On the low-radial arc the other endpoint is $(t_e,1-t_e)$, giving
$\lambda_{\mathrm{lo}}(e)$.  With $c=1-e$ and
$t_e=2c/(1+\sqrt{16c^2-3})$, the inequality
$\lambda_{\mathrm{lo}}(e)>1/3$ is equivalent to
$4-c>7t_e$.  After isolating and squaring its positive radical, the
difference factors as
\[
 4(2c-1)(2c^3-15c^2-4c+16)>0
 \qquad\left(\frac12<c\le\frac{\sqrt3}{2}\right);
\]
the cubic decreases on this interval and at $c=\sqrt3/2$ equals
$(19-5\sqrt3)/4>0$.  Finally the low-radial range gives
$e\ge1-\sqrt3/2>2/15$, and therefore
$1/3>1-5e$ whenever this comparison is needed.
\end{proof}

\paragraph{Body consequence.}
The two strict chord estimates supply the successive return bounds in the
Appendix D certificate for \zcref{prop:readable-k410-ce1}.

\subsection{Module D.6: Vd corner margins}

\paragraph{Exact interface.}
Use the Vd1/Vd2 corner-form actual reaches and the radial separation objects
of \zcref{prop:finite-vd-radial-separation}; selected lower bounds remain
distinguished from the actual reaches $A,B$.

\paragraph{Variables and feasible set.}
Take $0<A\le1$ and $0\le B\le1$ from the actual corner traces and
$\theta>0$, with $\Delta_\theta=\sqrt{\theta^2+\theta+1}$.  The endpoint
$u_+$ is used only when $\Haus(T_0\cap r_1)>0$.  When selected demands are
used, write them in lowercase and require $A\ge a_*$ and $B\ge b_*$.

\paragraph{Objective.}
Evaluate the corner's own-radial reach and far supported endpoint, and prove
the strict margins needed to place the residual radial point beyond both.

\paragraph{Exact result.}

\begin{lemma}[Vd corner margins]
\label{lem:app-vd-corner-margins}
In the Vd1/Vd2 corner form, let the actual boundary reaches satisfy
$0<A\le1$ and $0\le B\le1$, let $\theta>0$, and put
$\Delta_\theta=\sqrt{\theta^2+\theta+1}$.  The own-radial reach is
\[
 C_0=\frac{\Delta_\theta-A-\theta B}{\theta+1}.
\]
If $\Haus(T_0\cap r_1)>0$, the far endpoint of that positive forward
support is
\[
 u_+=\frac{\Delta_\theta-A-\theta B-1}{\theta}.
\]
Consequently
\[
 C_0<1-\min\{A,B\},
 \qquad u_+<1-B
 \quad\text{when }\Haus(T_0\cap r_1)>0.
\]
If $A\ge a_*$ and $B\ge b_*$, then
\[
 C_0<1-\min\{a_*,b_*\},
 \qquad u_+<1-b_*
 \quad\text{when }\Haus(T_0\cap r_1)>0.
\]
Reflection supplies the backward-support version.
\end{lemma}

\paragraph{Solution.}

\begin{proof}[Proof of \zcref{lem:app-vd-corner-margins}]
Since $\Delta_\theta<\theta+1$,
\[
 C_0<1-\frac{A+\theta B}{\theta+1}\le1-\min\{A,B\},
\]
while
\[
 u_+<1-B-\frac A\theta<1-B.
\]
Monotonicity gives the selected-demand version.
\end{proof}

\paragraph{Body consequence.}
The strict corner margins are used by the nonadjacent branch of
\zcref{prop:finite-vd-radial-separation}.
